% =====================================================================
\chapter{Técnicas de integración}
\label{tema:3}
% =====================================================================

La regla de Barrow reduce el cálculo de $\int_a^b f$ a encontrar una primitiva de $f$. A diferencia de la derivación, que sigue reglas mecánicas, la búsqueda de primitivas es un arte con unas pocas herramientas: el cambio de variable, que invierte la regla de la cadena; la integración por partes, que invierte la regla del producto; la descomposición en fracciones simples, que resuelve por completo las funciones racionales; y las identidades trigonométricas. Algunas funciones elementales, como $e^{-x^{2}}$, no tienen primitiva elemental, y para ellas se recurre a las funciones definidas por integrales del \cref{tema:2} o a las series del \cref{tema:6}. Se siguen el capítulo~19 de Spivak~\cite{spivak}, las §5.5 y~5.6 de~\cite{openstax1} y §3.1 a~3.5 de~\cite{openstax2}, el capítulo~7 de Strang~\cite{strang} y la unidad~4 del curso 18.01SC~\cite{mit1801}.

\section{Cambio de variable}

\begin{teorema}[Cambio de variable]\label{teo:sustitucion}
Sea $\varphi\colon[\alpha,\beta]\to\R$ con derivada continua, y sea $f$ continua en un intervalo $I$ que contiene a $\varphi([\alpha,\beta])$. Entonces
\[
\int_{\alpha}^{\beta}f\bigl(\varphi(t)\bigr)\,\varphi'(t)\dif t=\int_{\varphi(\alpha)}^{\varphi(\beta)}f(x)\dif x .
\]
En forma indefinida: si $F'=f$, entonces $\int f(\varphi(t))\,\varphi'(t)\dif t=F(\varphi(t))+C$.
\end{teorema}

En la práctica se escribe $u=\varphi(t)$, $\dif u=\varphi'(t)\dif t$, y se sustituye formalmente. En una integral definida conviene cambiar también los límites, lo que evita deshacer el cambio.

\addtocounter{definicion}{1}%

El teorema también se usa en sentido inverso: para calcular $\int_a^b f(x)\dif x$ se escribe $x=\varphi(t)$ con $\varphi$ biyectiva de $[\alpha,\beta]$ en $[a,b]$, $\varphi(\alpha)=a$, $\varphi(\beta)=b$, y se integra en $t$.

\addtocounter{definicion}{1}%

\section{Integración por partes}

\begin{teorema}[Integración por partes]\label{teo:partes}
Sean $u,v$ con derivada continua en $[a,b]$. Entonces
\[
\int_a^b u(x)\,v'(x)\dif x=\eval{u(x)\,v(x)}{a}{b}-\int_a^b u'(x)\,v(x)\dif x ,
\]
y en forma indefinida, $\int u\,v'=u\,v-\int u'\,v$, o abreviadamente $\int u\dif v=uv-\int v\dif u$.
\end{teorema}

La técnica es útil cuando $u$ se simplifica al derivarla (un polinomio, un logaritmo, una función trigonométrica inversa) y $v'$ no se complica al integrarla (una exponencial, un seno, un coseno).

\addtocounter{definicion}{1}%

La integración por partes repetida conduce a \emph{fórmulas de reducción}, que expresan una integral con un parámetro $n$ en función de la misma integral con $n-1$ o $n-2$ y se resuelven por inducción (\cref{eje:3-reduccion,eje:3-wallis}).

\section{Funciones racionales y fracciones simples}

Una función racional $P/Q$ es \emph{propia} si $\operatorname{gr}P<\operatorname{gr}Q$. Toda función racional es un polinomio más una función racional propia, por la división de polinomios. Todo polinomio real se factoriza en factores lineales $(x-r)$ y cuadráticos irreducibles $(x^{2}+px+q)$, con $p^{2}-4q<0$.

\begin{teorema}[Descomposición en fracciones simples]\label{teo:fracciones}
Sea $P/Q$ una función racional propia y supóngase $Q$ mónico factorizado como producto de potencias de factores lineales $(x-r)^{k}$ y cuadráticos irreducibles $(x^{2}+px+q)^{m}$. Entonces $P/Q$ es, de forma única, una suma de términos
\[
\frac{A_j}{(x-r)^{j}}\ \ (1\le j\le k)\qquad\text{y}\qquad\frac{B_j\,x+C_j}{(x^{2}+px+q)^{j}}\ \ (1\le j\le m),
\]
uno por cada potencia de cada factor.
\end{teorema}

La demostración general es algebraica~\cite[cap.~19]{spivak}. El caso más frecuente, con raíces reales simples, admite una prueba breve que además da los coeficientes.

\begin{proposicion}[Raíces simples]\label{prop:raices-simples}
Sea $Q(x)=(x-r_1)\cdots(x-r_n)$ con $r_1,\dots,r_n$ distintos y sea $P$ un polinomio de grado menor que $n$. Entonces
\[
\frac{P(x)}{Q(x)}=\sum_{i=1}^{n}\frac{A_i}{x-r_i},\qquad A_i=\frac{P(r_i)}{\prod_{j\neq i}(r_i-r_j)}=\frac{P(r_i)}{Q'(r_i)} .
\]
\end{proposicion}

Las fracciones simples tienen primitivas conocidas: $\int\frac{A\dif x}{x-r}=A\ln\abs{x-r}+C$; $\int\frac{A\dif x}{(x-r)^{j}}=-\frac{A}{(j-1)(x-r)^{j-1}}+C$ si $j\ge2$; y para el factor cuadrático se completa el cuadrado, $x^{2}+px+q=(x+\frac p2)^{2}+s^{2}$ con $s^{2}=q-\frac{p^{2}}{4}>0$, y se separa la parte que es derivada del denominador:
\[
\int\frac{\dif x}{(x+h)^{2}+s^{2}}=\frac1s\arctan\frac{x+h}{s}+C,\qquad\int\frac{2x+p}{x^{2}+px+q}\dif x=\ln(x^{2}+px+q)+C .
\]

\addtocounter{definicion}{1}%

\section{Integrales trigonométricas}

Para $\int\sin^{m}x\cos^{n}x\dif x$ con $m,n\in\N_0$ hay dos casos~\cite[§3.2]{openstax2}.
\begin{enumerate}
  \item Si alguno de los exponentes es impar, por ejemplo $m=2k+1$, se aparta un factor $\sin x$, se escribe $\sin^{2k}x=(1-\cos^{2}x)^{k}$ y se sustituye $u=\cos x$. El caso $n$ impar es simétrico, con $u=\sin x$.
  \item Si ambos son pares, se reduce el grado con $\sin^{2}x=\frac{1-\cos2x}{2}$, $\cos^{2}x=\frac{1+\cos2x}{2}$ y $\sin x\cos x=\frac{\sin2x}{2}$.
\end{enumerate}

\addtocounter{definicion}{1}%

Las raíces cuadráticas se eliminan con \emph{sustituciones trigonométricas}~\cite[§3.3]{openstax2}: $\sqrt{a^{2}-x^{2}}$ con $x=a\sin t$, como en el \cref{ej:circulo}; $\sqrt{a^{2}+x^{2}}$ con $x=a\tan t$; y $\sqrt{x^{2}-a^{2}}$ con $x=a/\cos t$. Por último, la sustitución $t=\tan\frac x2$ transforma cualquier función racional de $\sin x$ y $\cos x$ en una función racional de $t$, pues $\sin x=\frac{2t}{1+t^{2}}$, $\cos x=\frac{1-t^{2}}{1+t^{2}}$ y $\dif x=\frac{2\dif t}{1+t^{2}}$.

\section{Estrategia}

Ante una integral conviene probar, por este orden: (1) simplificar algebraicamente y buscar una primitiva inmediata (\cref{tab:inmediatas}); (2) buscar una función y su derivada en el integrando, para un cambio de variable; (3) si es un producto de un polinomio o logaritmo por una exponencial o función trigonométrica, integrar por partes; (4) si es racional, dividir y descomponer en fracciones simples; (5) si es trigonométrica o contiene raíces de polinomios de grado~2, usar identidades o sustituciones trigonométricas. Toda primitiva se comprueba derivando.

\addtocounter{definicion}{1}%

\addtocounter{definicion}{1}%

